%! Author = lili (Transkription)
%! Date = 10/7/26

\chapter{Testen von Hypothesen}\label{ch:testen-von-hypothesen}
\untit{Vorlesung 10. Juli 2026}

\txc{(2026)} \textbf{Ziel:} Verwende Daten, um zu entscheiden, ob eine zuvor formulierte Annahme über die
\gls{verteilung} bestätigt werden kann.
Ein \textbf{statistischer Test} ist ein Verfahren dafür.

\section{Abstrakte Formulierung für das Vorgehen}\label{sec:abstrakte-formulierung-hypothesentest}

\begin{itemize}
    \item Ereignisraum $\eraum$
    \item Stichprobenraum $\spr$ (bei Stichprobenlänge $n$)
    \item Parameterraum $\praum$
    \item statistisches Modell $\pefam = \{\pe_{\param} : \param\in\praum\}$, gegeben durch Wahrscheinlichkeitsfunktion
    $\wfk_{\param}$ oder Dichte $f_{\param}$
    \item Stichprobenvariable $X_1,\dots,X_n : \eraum \to \spr$, $X(\ele) = (X_1(\ele),\dots,X_n(\ele))$
    \item Stichprobe $x = (x_1,\dots,x_n) \in \spr$
    \item Hypothese: $\param \in H_0$
    \item Alternative: $\param \in H_A$, mit $H_0 \cap H_A = \emptyset$ ($H_0 \cup H_A = \praum$ nicht notwendig)
    \item Teststatistik
    \item Test
    \item Entscheidungsregel
\end{itemize}

\begin{bsp}
    \textbf{Mensa} \txc{(2026)} \\
    20\% der Studierenden sind mit dem alten Mensa-Konzept unzufrieden.
    Es wird ein neues Konzept eingeführt.
    Nach einer Woche werden 20 Studierende befragt; 8 davon sind unzufrieden.

    \paragraph{Frage.} Wird das neue Angebot als schlechter empfunden -- oder ist unsere Stichprobe nicht
    repräsentativ?

    \paragraph{Modell.}
    \[
        \eraum = \{0,1\}^{n}, \qquad n = 20,
    \]
    ($\ele_i=1$: $i$-te Person unzufrieden).
    \[
        \praum = [0,1], \qquad \wfk_{\param}(\ele) = \param^{\sum_{i=1}^n \ele_i} (1-\param)^{n-\sum_{i=1}^n
        \ele_i} \cdot
    \]
    Stichprobenvariable $X_i$: Bernoulli-Variable, unter $\param$ gilt $\pe_{\param}(X_i=1)=\param$
    ($i$-te Person unzufrieden).
\end{bsp}

\karkar{bem_ziel_statistischer_test}
\karkar{bem_bestandteile_testtheorie_vokabular}

\section{Wahl der Hypothesen}\label{sec:wahl-der-hypothesen}

\txc{(2026)} So wählen wir die Hypothesen, dass die \textbf{Nichtablehnung} (Nicht-Verwerfen) das Ziel ist.

Die Mensabetreiber wollen sich daran erfreuen, dass die Studierenden mit dem neuen Konzept zufriedener sind.
Daher
ist das Ziel, zu widerlegen, dass der Anteil Unzufriedener gleich geblieben ist:
\[
    H_0 : \param = 0.2 \quad \text{(Unzufriedenanteil unverändert)}
\]
\[
    H_A : \param < 0.2 \quad \text{(kleiner)}
\]

Skeptische Vertreter der Studierenden dagegen würden anders vorgehen.
Sie wollen ebenfalls widerlegen, dass der
Anteil Unzufriedener gleich geblieben ist -- aber sie werden den Vergleich mit einer \emph{größeren} Zahl
Unzufriedener ziehen wollen:
\[
    H_0 : \param = 0.2 \quad \text{(unverändert)}
\]
\[
    H_A : \param > 0.2 \quad \text{(größer)}
\]

\begin{defi}
    \textbf{Neyman--Pearson-Formulierung der Testtheorie} \label{defi:neyman-pearson} \txc{(2026)} \\
    Sei $\pefam = \{\pe_{\param} : \param \in \praum\}$ ein statistisches Modell.
    Gegeben $H_0, H_A \subseteq \praum$
    mit $H_0 \cap H_A = \emptyset$.
    \begin{itemize}
        \item Hypothese: $\param \in H_0$
        \item Alternative: $\param \in H_A$
    \end{itemize}
    Das Entscheidungsproblem $H_0$ gegen $H_A$ heißt \textbf{Testproblem}. $H_0$ heißt \textbf{Nullhypothese}.
\end{defi}

\paragraph{Sprechweise.} Ist $H_0$ einelementig wie in unserem Mensa-Beispiel, so heißt die Hypothese
\textbf{einfach}.
Andernfalls heißt sie \textbf{zusammengesetzt}.
Wir sprechen genauso von einer einfachen oder
einer zusammengesetzten Alternative.
Im Mensa-Beispiel ist die Alternative zusammengesetzt.

Grundsätzlich wollen wir die Nullhypothese \textbf{konservativ} wählen: Sie ist in der Regel das Gegenteil von dem
vermuteten Zusammenhang, den wir aufgrund der Beobachtungen rechtfertigen wollen.
Wird nun in unserem
Hypothesentest die Nullhypothese abgelehnt, so können wir den Zusammenhang als \textbf{statistisch signifikant}
betrachten.

\karkar{bem_gleiche_h0_verschiedene_ha}
\karkar{def_testproblem_neyman_pearson}
\karkar{bem_muss_h0_und_ha_ganz_theta_abdecken}
\karkar{def_einfache_zusammengesetzte_hypothese}
\karkar{bem_wahl_nullhypothese_konservativ}

\section{Teststatistik}\label{sec:teststatistik}

\txc{(2026)} Analog zum Punktschätzer führen wir nun eine Abbildung ein, die die vorhandenen Daten in eine Zahl
verwandelt.

\begin{defi}
    \textbf{Teststatistik} \txc{(2026)} \\
    Eine Abbildung $T : \spr \to \rz$ heißt \textbf{Teststatistik}.
    Wähle $T$ so, dass \emph{große} Werte von
    $T(x)$ dafür sprechen, dass wir eine unter $H_0$ \textbf{atypische} Stichprobe sehen.
\end{defi}

\begin{bsp}
    \textbf{Mensa (Fortsetzung)} \txc{(2026)} \\
    Viele Unzufriedene in der Stichprobe $x$ bedeuten insbesondere, dass $\bar x = \frac1n\sum_{i=1}^n x_i$ groß
    ist.

    Beim Vergleich von $\param=0.2$ (unter $H_0$) zu $\param \in H_A$, d.\,h.\ $\param < 0.2$ (Mensabetreiber),
    sind große Werte von $\bar x$ für $\param=0.2$ eher untypisch für die Alternative.
    Daher ist $\spmw$ keine
    geeignete Teststatistik für den Test $H_0$ gegen $H_A$ in diesem Fall.
    Wir wählen stattdessen
    \[
        \widetilde T = 1 - \spmw \cdot
    \]
    Die Wahl $\widetilde T$ entspricht einem Tausch von $0$ für Zufriedene und $1$ für Unzufriedene in der
    Interpretation von $\ele \in \{0,1\}^n$, hin zu $1$ für Zufriedene und $0$ für Unzufriedene.

    Die Studierendenvertreter dagegen wählen $T = \spmw$, weil im Vergleich zwischen $\param=0.2$ (unter $H_0$)
    und $\widetilde\param \in H_A$, d.\,h.\ $\param > 0.2$, große Werte von $T$ gegen die Nullhypothese sprechen.
\end{bsp}

\begin{defi}
    \textbf{$p$-Wert} \label{defi:p-wert} \txc{(2026)} \\
    \[
        p(x) = \sup_{\param \in H_0} \pe_{\param}\big( T(X) \geq T(x) \big) \cdot
    \]
    Der $p$-Wert drückt aus, wie wahrscheinlich es unter der Nullhypothese \emph{maximal} ist, dass die
    theoretische Teststatistik $T(X)$ mindestens so groß ist wie die Teststatistik, ausgewertet in der Stichprobe
    $x$.
    So wie wir die Teststatistik wählen, erwarten wir einen \textbf{kleinen} $p$-Wert.
\end{defi}

\begin{bsp}
    \textbf{Mensa (Fortsetzung)} \txc{(2026)} \\
    Für den Mensabetreiber-Test ergab sich ein $p$-Wert nahe $1$ -- das ist absurd groß, sollte uns aber nicht
    wundern, denn statt weniger sind in der Stichprobe \emph{mehr} Unzufriedene als unter $H_0$ erwartet.

    Für den Studierendenvertreter-Test ergab sich dagegen ein kleiner $p$-Wert: Unter der Nullhypothese ist es im
    Vergleich zur Alternativhypothese sehr unwahrscheinlich, so viele Unzufriedene anzutreffen.
    Dies spricht dafür,
    dass die Nullhypothese abgelehnt werden sollte.
    Beweist also nicht die Gültigkeit der Alternative.
\end{bsp}
\komm{
    Die konkreten $p$-Werte selbst waren in der vorliegenden Quelle nicht zuverlässig lesbar und wurden daher hier
    nicht wiedergegeben.
}

\karkar{def_teststatistik_th}
\karkar{bem_wahl_teststatistik_haengt_von_alternative_ab}
\karkar{def_p_wert_th}
\karkar{bem_p_wert_intuition_th}

\section{Test, Signifikanzniveau, Verwerfungsbereich}\label{sec:test-signifikanzniveau-verwerfungsbereich}

\begin{defi}
    \textbf{Test} \txc{(2026)} \\
    Ein \textbf{Test} ist eine Abbildung $\varphi : \spr \to \{0,1\}$. $\varphi(x)=1$ bedeutet: Hypothese $H_0$ ist
    zu verwerfen. $\varphi(x)=0$ bedeutet: $H_0$ wird nicht verworfen.
\end{defi}

\begin{defi}
    \textbf{Signifikanzniveau} \txc{(2026)} \\
    Das Signifikanzniveau eines Tests $\varphi$ ist die Zahl
    \[
        \sup_{\param \in H_0} \pe_{\param}\big( \varphi(X) = 1 \big),
    \]
    d.\,h.\ die maximale Wahrscheinlichkeit, dass der Test unter $H_0$ verwirft.
\end{defi}

\begin{bsp}
    \textbf{Mensa (Fortsetzung)} \txc{(2026)} \\
    Für den Mensabetreiber-Test ist $\varphi$ per Konstruktion ein Test zum Niveau $\alpha \in (0,1)$.
    Für
    $\alpha=0.05$ verwirft der Test $H_0$ \emph{nicht}: Die Mensabetreiber können nicht behaupten, dass sich die
    Zufriedenheit verbessert hat.
    Alles andere wäre auch absurd.

    Für den Studierendenvertreter-Test verwirft $\varphi$ die Nullhypothese zum Niveau $\alpha=0.05$ -- und
    tatsächlich für jedes hinreichend große Niveau.
    Das neue Konzept der Mensa darf als signifikant angesehen
    werden für die gewachsene Unzufriedenheit.
\end{bsp}

Der $p$-Wert enthält mehr Information als das Niveau eines Tests.
Gute Statistik gibt den $p$-Wert an.

\begin{bem}
    \textbf{Verwerfungsbereich} \txc{(2026)} \\
    \[
        C_\alpha = \{ x \in \spr : \varphi(x)=1 \}
    \]
    ist die Menge aller Stichproben, für die der Test $\varphi$ die Nullhypothese verwirft.

    Wir könnten daher auch $C_\alpha$ bestimmen und anschließend prüfen, ob unsere Stichprobe in $C_\alpha$ liegt.
\end{bem}

\begin{bsp}
    \textbf{Mensa (Fortsetzung), $\alpha=0.05$} \txc{(2026)} \\
    Hätte die Stichprobe eine (in der Quelle nicht mehr sicher lesbare) kleinere Anzahl Unzufriedener ergeben, so
    hätten auch die Studierendenvertreter ihre Hypothese $H_0$ nicht verwerfen können.
\end{bsp}

\karkar{def_test_abbildung}
\karkar{def_signifikanzniveau_th}
\karkar{bem_p_wert_mehr_information_als_niveau}
\karkar{def_verwerfungsbereich_th}

\begin{defi}
    \textbf{Fehler 1.\ und 2.\ Art} \txc{(2026)} \\
    \begin{itemize}
        \item \textbf{Fehler 1.\ Art:} Eine gültige Hypothese wird verworfen.
        D.\,h.\ $\param \in H_0$ und
        $\varphi(X)=1$.
        \item \textbf{Fehler 2.\ Art:} Eine ungültige Hypothese wird nicht verworfen.
        D.\,h.\ $\param \in H_A$ und
        $\varphi(X)=0$.
    \end{itemize}
\end{defi}

\begin{bem}
    \txc{(2026)} \\
    Die Wahrscheinlichkeit für einen Fehler 1.\ Art ist beschränkt durch das Signifikanzniveau des Tests.
\end{bem}

\paragraph{Ziel.} Finde, zu gegebenem Signifikanzniveau $\alpha$, einen Test derart, dass sein
Signifikanzniveau den vorgegebenen Wert nicht überschreitet.
Idealerweise ist auch die Wahrscheinlichkeit für einen
Fehler 2.\ Art nicht zu groß.

\begin{bsp}
    \textbf{Mensa, Studierendenvertreter} \txc{(2026)} \\
    Test zum Niveau $\alpha=0.05$.
    Wir könnten stattdessen eine andere Entscheidungsregel wählen, um die
    Wahrscheinlichkeit für einen Fehler 2.\ Art zu reduzieren -- allerdings nur um den Preis, dass die
    Wahrscheinlichkeit für einen Fehler 1.\ Art (bzw.\ das Signifikanzniveau) größer wird.
\end{bsp}
\komm{
    Die Vorlesung berechnete an dieser Stelle die konkrete Fehlerwahrscheinlichkeit 1.\ Art sowie die
    Fehlerwahrscheinlichkeit 2.\ Art als Funktion von $\param$ (letztere geht gegen $0$ für $\param\to1$ und gegen
    $1-\alpha$ für $\param\to0.2$). Die genauen Zwischenwerte waren in der Quelle nicht zuverlässig lesbar.
}

\karkar{def_fehler_1_art}
\karkar{def_fehler_2_art}
\karkar{bem_fehler_1_art_beschraenkt_durch_niveau}
\karkar{bem_ziel_testkonstruktion}
\karkar{bem_tradeoff_fehler_1_2_art}

\section{Vorgehen bei einem Hypothesentest}\label{sec:vorgehen-hypothesentest}

\txc{(2026)}
\begin{enumerate}
    \item statistisches Modell / Verteilungsannahme
    \item Hypothese und Alternative formulieren
    \item Signifikanzniveau wählen
    \item Teststatistik konstruieren
    \item Entscheidungsregel festlegen
    \item Test zum gewählten Signifikanzniveau, Stichprobe einsetzen
    \item Schlussfolgerung und $p$-Wert angeben
\end{enumerate}

\begin{bem}
    \textbf{Asymmetrie} \txc{(2026)} \\
    \begin{itemize}
        \item $\varphi(x)=1$: $H_0$ verwerfen.
        \item $\varphi(x)=0$: $H_0$ nicht verwerfen -- \textbf{keine Aussage über die Alternative}.
    \end{itemize}
    Die Entscheidung \emph{für} eine Hypothese wäre weniger sicher als die Entscheidung \emph{gegen} eine
    Hypothese.
    ``Nicht verwerfen'' bedeutet nur, dass die Stichprobe es nicht erlaubt, die Hypothese mit
    hinreichender Sicherheit zu verwerfen.
\end{bem}

\karkar{bem_rezept_hypothesentest}
\karkar{bem_asymmetrie_testentscheidung}

\section*{Begleitbeispiel: Geburtenstatistik}\label{sec:begleitbeispiel-geburtenstatistik}

Das folgende Beispiel soll uns als Begleitung an die allgemeine Testtheorie dienen, nachdem wir am Mensa-Beispiel
die Grundlage der Testtheorie gesehen haben.

\begin{bsp}
    \textbf{Geburtenstatistik} \txc{(2026)} \\
    Geburtenstatistik aus der Schweiz, 1950--1970.
    Naiverweise wird oft angenommen, dass die Hälfte aller
    geborenen Kinder Jungen sind.
    Dies sei unsere Nullhypothese.
    Glauben die Daten, diese Hypothese zu verwerfen?
    Signifikanzniveau $\alpha=0.01$.
    \[
        H_0 : \param = \frac12 \qquad \text{gegen} \qquad H_A : \param \ne \frac12 \cdot
    \]

    Unter $H_0$ gilt für jedes geborene Kind: $\pe(\text{Junge}) = \param = \frac12$. $X_i$ ($i=1,\dots,n$) mit
    $X_i=1$, falls das $i$-te Kind ein Junge ist, ist eine Bernoulli-Variable mit $\param = \E{X_i}$,
    \[
        \Var(X_i) = \param(1-\param) \cdot
    \]
    $S_n = \sum_{i=1}^n X_i$: Anzahl der geborenen Jungen.

    \paragraph{Wir suchen eine Teststatistik, bei der große Werte unter $H_0$ atypisch sind.}

    \textit{Idee:} Unter $H_0$ sollte $\left|\frac1n S_n - \param_0\right|$ klein sein, weil nach dem Gesetz der
    großen Zahlen $\frac1n S_n$ für großes $n$ nahe an $\param$ sein wird.
    Außerdem wollen wir vom zentralen
    Grenzwertsatz profitieren:
    \[
        \sqrt{\frac{n}{\param_0(1-\param_0)}} \left( \frac1n S_n - \param_0 \right)
    \]
    ist für große $n$ näherungsweise standardnormalverteilt.
    Wir wählen daher
    \[
        T(x) = \sqrt{\frac{n}{\param_0(1-\param_0)}}\ \left| \frac{S_n}{n} - \param_0 \right|
    \]
    als Teststatistik.

    \paragraph{Test zum Signifikanzniveau $\alpha=0.01$.}
    \[
        \varphi(x) = \begin{cases}
                         1 & \text{falls } T(x) \geq z_{1-\alpha/2} \\
                         0 & \text{sonst}
        \end{cases}
    \]
    ($z_{1-\alpha/2}$: $(1-\alpha/2)$-Quantil der $N(0,1)$-Verteilung).
    Verwerfungsbereich: $|T| > z_{1-\alpha/2}$.

    Unter $H_0$ ($\param=\frac12$) ergab die Auswertung für die tatsächliche Stichprobe $\varphi(x)=1$: $H_0$
    muss verworfen werden.

    Der in diesem Beispiel durchgeführte Test ist ein \textbf{zweiseitiger Gauß-Test}.
\end{bsp}
\komm{
    Die genauen Zahlenwerte (Stichprobengröße, Anzahl Jungen, konkreter Wert von $T(x)$) waren in der Quelle
    nicht zuverlässig lesbar und wurden daher nicht übernommen.
}

Wir unterscheiden:
\begin{itemize}
    \item \textbf{einseitige} Tests wie im Mensa-Beispiel, mit $\param\in H_0$ und $\param\in H_A \Leftrightarrow
    \param < \param_0$ (linksseitig) bzw.\ $\param \in H_A \Leftrightarrow \param > \param_0$ (rechtsseitig);
    \item \textbf{zweiseitige} Tests wie im Geburten-Beispiel, wenn $\param\in H_0$ und $\param \ne \param_0$.
\end{itemize}

Im Folgenden befassen wir uns mit Tests, die Hypothesen zum Erwartungswert betreffen.
Seien $X_1,\dots,X_n$
u.i.v.\ Stichprobenvariable.

\karkar{bem_konstruktion_teststatistik_ggz_zgws}
%\karkar{bem_geburtenbeispiel_typ}
\karkar{bem_einseitig_vs_zweiseitig_unterscheidung}

\section{Einfacher Gauß-Test}\label{sec:einfacher-gausstest}

Anwendbar für $X_1,\dots,X_n$ normalverteilt mit unter $H_0$ bekannter Varianz $\sigma^2$.

Wir wenden diesen Test wie im Geburten-Beispiel auch bei Bernoulli-verteilten $X_1,\dots,X_n$ an, bei großem $n$.
Unter $H_0$ kennen wir $\param_0$ und damit auch die Varianz $\sigma^2 = \param_0(1-\param_0)$.

\subsection*{a) Zweiseitiger Test}

\[
    H_0 : \wfk = \wfk_0 \qquad \text{gegen} \qquad H_A : \wfk \ne \wfk_0 \cdot
\]
Teststatistik wie im Geburten-Beispiel:
\[
    T(x) = \sqrt n\, \frac{| \bar x - \wfk_0 |}{\sigma}
\]
$T(X)$ ist bei normalverteilten Stichprobenvariablen der Betrag einer standardnormalverteilten \gls{zv}.
Bei
Bernoulli-verteilter Stichprobenvariable ist $T$ nur \emph{näherungsweise} der Betrag einer
standardnormalverteilten \gls{zv}.

Für den Test $\varphi$ zum Signifikanzniveau $\alpha$ wählen wir
\[
    \varphi(x) = \begin{cases} 1 & \text{falls } T(x) \geq z_{1-\alpha/2} \\ 0 & \text{sonst} \end{cases}
\]
Verwerfungsbereich: $|T| > z_{1-\alpha/2}$.

\komm{
Man könnte auch schreiben:
Teststatistik wie im Geburten-Beispiel:
    \[
        T(x) = \sqrt n\, \frac{ \bar x - \wfk_0 }{\sigma}
    \] und dann
    \[
        \varphi(x) = \begin{cases} 1 & \text{falls } |T(x)| \geq z_{1-\alpha/2} \\ 0 & \text{sonst} \end{cases}
    \]
}
\subsection*{b) Einseitiger Test -- Anpassungen}

Im Fall $\param \leq H_0\ (\param_0)$, $\param > H_A$
\txc{(Das ist rechtsseitig, siehe $\param > H_A \Leftrightarrow H_A < \param$)}:
\[
    \varphi(x) = \begin{cases} 1 & \text{falls } T(x) \geq z_{1-\alpha} \\ 0 & \text{sonst} \end{cases}
    \qquad \text{mit } T(x) = \sqrt n\, \frac{\bar x - \param_0}{\sigma} \cdot
\]
Verwerfungsbereich: $T > z_{1-\alpha}$.

\komm{
Merken: Wenn  $H_A: \param > \param_0$ ist, also $H_A: \param_0 < \param$, das Theta also auf der RECHTEN Seite steht, wenn man von klein zu groß geht, dann
ist es rechtsseitig. \\
In dem Fall sollte (wenn man von klein zu groß geht) auch $T(x)$ bzw. $T$ rechts stehen. \\ Analog für linksseitig...
}

Im Fall $\param \geq H_0$, $\param < H_A$:
\[
    \varphi(x) = \begin{cases} 1 & \text{falls } T(x) \leq -z_{1-\alpha} \\ 0 & \text{sonst} \end{cases}
\]
Verwerfungsbereich: $T < -z_{1-\alpha}$.
\komm{
$T(x) < -z_{1-\alpha} \Leftrightarrow T(x) < z_{\alpha}$
}

\komm{
    Übersicht:
    \[
        \begin{array}{lll}
            H_A:\mu>\mu_0
            &\Rightarrow&
            T>z_{1-\alpha}
            \\
            H_A:\mu<\mu_0
            &\Rightarrow&
            T<-z_{1-\alpha}
            \\
            H_A:\mu\neq\mu_0
            &\Rightarrow&
            |T|>z_{1-\alpha/2}
        \end{array}
    \]
}

\karkar{bem_gauss_test_voraussetzung}
\karkar{bem_gauss_test_bei_bernoulli}
\karkar{formel_teststatistik_gauss_test}
\karkar{bem_verteilung_teststatistik_gauss_normal_vs_bernoulli}
\karkar{formel_gauss_test_zweiseitig}
\karkar{formel_gauss_test_einseitig}

\section{Einfacher $t$-Test}\label{sec:einfacher-t-test}

Anwendbar für $X_1,\dots,X_n$ normalverteilt mit unbekannter Varianz $\sigma^2$.

Das Vorgehen ist analog zum einfachen Gauß-Test.
Nun müssen wir in der Teststatistik die Varianz durch ihren
Schätzer $\widehat\sigma^2 = \dfrac{1}{n-1}\displaystyle\sum_{i=1}^n (x_i-\bar x)^2$ ersetzen und folglich mit den
Quantilen der $t$-Verteilung arbeiten.

\subsection*{a) Zweiseitiger Test}

\[
    H_0 : \wfk = \wfk_0 \qquad \text{gegen} \qquad H_A : \wfk \ne \wfk_0 \cdot
\]
Teststatistik:
\[
    T(x) = \sqrt n\, \frac{|\bar x - \wfk_0|}{\widehat\sigma} \cdot
\]
$T(X)$ ist der Betrag einer $t$-verteilten \gls{zv} mit $n-1$ Freiheitsgraden.

Für den Test zum Signifikanzniveau $\alpha$ wählen wir
\[
    \varphi(x) = \begin{cases} 1 & \text{falls } T(x) \geq t_{n-1,1-\alpha/2} \\ 0 & \text{sonst} \end{cases}
\]
\komm{
Auch hier könnte man schreiben
    \[
        T(x) = \sqrt n\, \frac{\bar x - \wfk_0}{\widehat\sigma} \cdot
    \] und dann
    \[
        \varphi(x) = \begin{cases} 1 & \text{falls } |T(x)| \geq t_{n-1,1-\alpha/2} \\ 0 & \text{sonst} \end{cases}
    \]
}


\subsection*{b) Einseitiger Test}

Analog zum Gauß-Test, mit $\sigma$ ersetzt durch $\widehat\sigma$ und $z_{1-\alpha}$ ersetzt durch
$t_{n-1,1-\alpha}$.

\komm{
    $T(x) < -t_{n-1,1-\alpha} \Leftrightarrow T(x) < t_{n-1,\alpha}$
}

\komm{
Übersicht:
    \[
        \begin{array}{lll}
            H_A:\mu>\mu_0
            &\Rightarrow&
            T>t_{n-1,1-\alpha}
            \\
            H_A:\mu<\mu_0
            &\Rightarrow&
            T<-t_{n-1,1-\alpha}
            \\
            H_A:\mu\neq\mu_0
            &\Rightarrow&
            |T|>t_{n-1,1-\alpha/2}
        \end{array}
    \]
}

\karkar{bem_t_test_voraussetzung}
\karkar{bem_t_test_vorgehen_analogie}
\karkar{formel_teststatistik_t_test}
\karkar{formel_t_test_entscheidungsregeln}

\section{Zwei-Stichprobentests}\label{sec:zwei-stichprobentests}

\subsection*{1.\ Fall: Separate Stichproben}

\begin{bsp}
    \txc{(2026)} \\
    Stellen Sie sich vor, wir messen bei $n$ Personen jeweils die Länge beider Beine.
    Uns interessiert, ob es
    einen systematischen Unterschied gibt.

    Wir haben 2 $n$-dimensionale Stichprobenvariable $(X_i,Y_i)$, wobei in unserem Beispiel $x_i$ die Länge des
    rechten und $y_i$ die Länge des linken Beins sei.

    Wir nehmen an, dass $(X_1,Y_1),\dots,(X_n,Y_n)$ unabhängig sind.
    Aber wir dürfen \emph{nicht} annehmen, dass
    $X_i$ und $Y_i$ (selbst, innerhalb einer Person) unabhängig sind.
\end{bsp}

$D_i := X_i - Y_i$ sind unabhängig, weil $(X_1,Y_1),\dots,(X_n,Y_n)$ unabhängig sind. $D_1,\dots,D_n$ u.i.v.,
$D_i$ normalverteilt (Annahme) mit Erwartungswert $\wfk$ und unbekannter Varianz.

\subsection*{a) Zweiseitiger Test}

\[
    H_0 : \wfk = 0 \qquad \text{gegen} \qquad H_A : \wfk \ne 0 \cdot
\]
Unter $H_0$ ist $D_i$ normalverteilt mit $\E{D_i}=0$ und unbekannter Varianz.
Wir wissen: $T = \sqrt n\,
\dfrac{\bar D}{\widehat\sigma}$ ist $t$-verteilt mit $n-1$ Freiheitsgraden.
Dabei haben wir
\[
    \bar D = \frac1n \sum_{i=1}^n (x_i-y_i), \qquad \widehat\sigma^2 = \frac{1}{n-1} \sum_{i=1}^n (d_i-\bar d)^2
\]
gesetzt.
\[
    \varphi(y,x) = \begin{cases}
                       1 & \text{falls } \left| \sqrt n\, \dfrac{\bar D}{\widehat\sigma} \right| \geq t_{n-1,1-\alpha/2} \\
                       0 & \text{sonst}
    \end{cases}
\]

\subsection*{b) Einseitige Tests}

Analog: für $\param \leq H_0\ (0)$, $\param > H_A$:
\[
    \varphi(x) = \begin{cases} 1 & \text{falls } T(x) \geq t_{n-1,1-\alpha} \\ 0 & \text{sonst} \end{cases}
    \qquad \text{mit } T(x) = \sqrt n\, \frac{\bar D}{\widehat\sigma} \cdot
\]
Für $\param \geq H_0$, $\param < H_A$:
\[
    \varphi(x) = \begin{cases} 1 & \text{falls } T(x) \leq -t_{n-1,1-\alpha} \\ 0 & \text{sonst} \end{cases}
\]

\karkar{bem_zwei_stichproben_situation}
\karkar{bem_differenzentrick}
\karkar{formel_gepaarter_t_test_zweiseitig}